\textbf{II -- UN METODO PER AVVANTAGGIARSI IN APERTURA: IL GAMBETTO}

Nella prima lezione abbiamo parlato del rapido sviluppo e del guadagno
di tempi in apertura. Un metodo comune per guadagnare tempi si ha
sacrificando materiale, di norma un pedone. Il sacrificio di pedone per
avvantaggiarsi nello sviluppo o guadagnare spazio sulla scacchiera si
chiama gambetto. Il termine fu usato per la prima volta da Ruy Lopez nel
trattato \emph{``Libro de la invencion liberal y arte del juego de
¥xedrez''} del 1561 e deriva dall'italiano (gamba, sgambetto).

Di solito è il Bianco a sacrificare un pedone per incrementare quel
lieve vantaggio che gli deriva dal diritto alla prima mossa ma sono
stati elaborati numerosi gambetti anche per il Nero.

Per capire la filosofia del gambetto si veda questo esempio:

\emph{Partita Italiana:} \textbf{1.e4 e5 2.¤f3 ¤c6 3.¥c4 ¥c5 4.c3} scopo
di questo tratto è di spingere successivamente in d4 per piazzare al
centro due pedoni.

\textbf{4... ¤f6 5.d4}

\bookdiagramplaceholder{assets/diagrams/image6.png}{6}

Il Bianco ha spinto un secondo pedone al centro. Sostenuto com'è dal
\textbf{§}c3 esso non potrà facilmente essere rimosso. Per ottenere
questo indubbio vantaggio ha però dovuto giocare due mosse di pedone in
apertura e il Nero ne ha approfittato per continuare lo sviluppo con
4... ¤f6.

Nel 1824, un lupo di mare, il capitano di marina William Davies Evans
ideò una continuazione che raggiungeva la stessa configurazione (c3, d4,
e4) senza che il Nero si fosse avvantaggiato nello sviluppo (4... ¤f6).
Naturalmente non si poteva raggiungere l'obiettivo senza dare nulla in
cambio; nacque così un interessante sacrificio di pedone:

\emph{Gambetto Evans:} \textbf{1.e4 e5 2.¤f3 ¤c6 3.¥c4 ¥c5 4.b4} ecco la
mossa di Evans; \textbf{4... ¥xb4 5.c3 ¥c5 6.d4}

\bookdiagramplaceholder{assets/diagrams/image7.png}{7}

Si confronti questa posizione con quella derivante dalla Partita
Italiana. Il Bianco non ha più il pedone b2 ma il Cavallo nero si trova
ancora nella casa di partenza (g8).

Naturalmente, se ora togliessimo tutti i pezzi dalla scacchiera il Nero
vincerebbe il finale in virtù del pedone in più ma \emph{``prima del
finale gli dei hanno creato il centro partita''} recita un celebre
aforisma di Tartakower.

È importante riflettere che negli scacchi il valore dei pezzi non è
assoluto ma solo indicativo: esso varia a seconda della loro
disposizione sulla scacchiera. Il controllo di una colonna o di una
diagonale possono valere quanto un pedone o più. Per intenderci un
Alfiere che controlla sette case vale di più di uno che ne controlla
solo due e un sacrificio è giustificato se serve a dare maggiore forza
ai pezzi rimanenti. È chiaro che alla fine quel che conta, a meno di non
dare matto, è il materiale che, prima o poi, deve essere recuperato, se
possibile con gli interessi, pena la perdita della partita.

Nel secolo scorso, quando le tecniche difensive non erano sofisticate
come lo sono oggi, i gambetti erano di gran moda; oggi hanno perso molto
del loro smalto ma non sono scomparsi dalla pratica dei tornei.
Piuttosto è cambiata l'idea che ne è alla base. Un tempo essi venivano
giocati per aggredire senza compromessi il Re avversario, oggi li si
gioca per ottenere vantaggi di o¢dine posizionale da far fruttare
a£dirittura in finale.

È ovvio che l'offerta di pedone può essere rifiutata. Chi l'accetta se
ne assume gli inevitabili rischi senza illudersi di arrivare facilmente
ad un finale vinto.

Se in risposta ad un gambetto il Nero gioca a sua volta un gambetto, si
parla di controgambetto. Alcuni manuali usano, a mio avviso meno
correttamente, il termine controgambetto per gambetto giocato dal Nero.

La rivista inglese ``Gambit'' ha classificato oltre trecento gambetti.
Oltre all'Evans, eccone altri tra i più giocati.

\textbf{1. Sacrifici operati dal Bianco}

\emph{Gambetto di Re:} \textbf{1.e4 e5 2.f4}

\bookdiagramplaceholder{assets/diagrams/image8.png}{8}

Il Gambetto di Re, considerato a buon diritto il Re dei gambetti perché
in voga da oltre cinque secoli è ancora oggi perfettamente giocabile. In
caso di accettazione del pedone (2... exf4), il Bianco ottiene di poter
piazzare un secondo pedone al centro (d2-d4) e di aprire la colonna
``f'' alla Torre. Dopo l'arrocco corto e la rimozione del §f4 il Bianco
potrà fare pressione su f7 che, essendo difeso solo dal Re, è il punto
più debole dello schieramento nero.

\textbf{2... exf4 3.h4} il Bianco mina il pedone f4 per sfruttare la
colonna f. Altre buone continuazioni sono: 3.¥c4 e 3.d4.

\emph{Gambetto Danese:} \textbf{1.e4 e5 2.d4 exd4 3.c3 £xc3 4.¥c4}

\bookdiagramplaceholder{assets/diagrams/image9.png}{9}

Il Bianco sacrifica addirittura due pedoni pur di aprire il maggior
numero di linee.

Difficile dimostrare la correttezza di questi sacrifici ma al Nero è
richiesto un gioco attento ed ene¢gico.

Vediamo un possibile seguito:

\textbf{4... ¤xb2 5.¥xb2 d5!} l'idea del Nero, con questa forte reazione
centrale, è quella di restituire entrambi i pedoni per superare le
difficoltà d'apertura.

\textbf{6.¥xd5 ¤f6 7.¥xf7+ ¢xf7 8.£xd8 ¥b4+ 9.£d2 ¥xd2+ 10.¤xd2 ¦e8} e
il Nero non ha più nulla da temere.

\emph{Gambetto Scozzese:} \textbf{1.e4 e5 2.¤f3 ¤c6 3.d4 exd4 4.¥c4}

\bookdiagramplaceholder{assets/diagrams/image10.png}{10}

Questo gambetto, già citato nei manoscritti di Gioacchino Greco (XVII
secolo) è stato in anni recenti arricchito di nuove idee e lo si può
incontrare anche in tornei di alto livello.

\emph{Gambetto Siciliano:} \textbf{1.e4 c5 2.b4}

\bookdiagramplaceholder{assets/diagrams/image11.png}{11}

Questa mossa, già citata in un manoscritto del 1623, è giocata solo
sporadicamente a livello magistrale. Sergio Mariotti ne ha fatto ogni
tanto uso per sorprendere gli avversari. Ecco un possibile seguito:

\textbf{2... ¤xb4 3.a3 bxa3?!} è meglio 3... d5!, una mossa buona per
controbattere molti gambetti.

\textbf{4.¤xa3 d6 5.¥b2 ¤f6 6.¥c4 ¤c6 7.£e2 g6 8.f3 ¥g7 9.d4 0-0 10.0-0}
e il Bianco ha compenso per il pedone sacrificato.

\emph{Gambetto Morra:} \textbf{1.e4 c5 2.d4}

\bookdiagramplaceholder{assets/diagrams/image12.png}{12}

Questo gambetto conduce a scontri violenti in cui non è agevole per il
Nero tenere sotto controllo la situazione.

\textbf{2... ¤xd4 3.c3 £xc3} l'offerta può essere rifiutata senza
svantaggi sia con 3... d5 che con 3... ¤f6.

\textbf{4.¤xc3 ¤c6 5.¤f3 e6 6.¥c4 d6 7.0-0 ¤f6 8.£e2 ¥e7 9.¦d1} con
iniziativa.

\textbf{2. Sacrifici operati dal Nero}

\emph{Gambetto Albin:} \textbf{1.d4 d5 2.c4 e5}

\bookdiagramplaceholder{assets/diagrams/image13.png}{13}

Tentativo violento, ma non del tutto corretto, di conquistare
l'iniziativa.

\textbf{3.£xe5 d4 4.¤f3} sarebbe un grave errore 4.e3? ¥b4+! 5.¥d2 £xe3
6.fxe3 \emph{(6.¥xb4? exf2+ 7.¢e2 fxg1=}\textbf{¤}\emph{+! 8.¢e1 £h4+
9.¢d2 ¤c6 10.¥c3 ¥g4 e vince)} 6... £h4+ 7.g3 £e4 8.£f3 ¥xd2+ 9.¤xd2
£xe5 con vantaggio del Nero.

\textbf{4... ¤c6 5.g3} preparandosi a piazzare l'Alfiere sulla grande
diagonale per fare pressione sul futuro arrocco lungo del Nero.

\textbf{5... ¥g4 6.¥g2 £d7 7.0-0 0-0-0} una mossa importante: prepara la
futura apertura della colonna h con le spinte h5-h4.

Però nello scontro violento che questa posizione configura la pratica ha
dimostrato che in genere l'attacco del Bianco arriva prima.

\emph{Gambetto di Budapest:} \textbf{1.d4 ¤f6 2.c4 e5}

\bookdiagramplaceholder{assets/diagrams/image14.png}{14}

Pur non essendo del tutto giustificata questa coraggiosa spinta di
pedone, anche in tempi recenti, ha fatto registrare alcune belle
vittorie.

\textbf{3.£xe5 ¤g4} un'altra possibilità è 3... ¤e4 (Variante
Fajarowicz) in cui il Nero rinuncia alla pressione su d5 per puntare di
più sull'attacco.

\textbf{4.¥f4 ¤c6 5.¤f3 ¥b4+ 6.¤bd2 £e7 7.a3 ¤gxe5 8.¤xe5} non 8.axb4??
per 8... ¤d3 matto!

\textbf{8... ¤xe5 9.e3 ¥xd2+ 10.£xd2 d6} e il Bianco può contare sulla
coppia degli Alfieri.

\emph{Gambetto From:} \textbf{1.f4 e5}

\bookdiagramplaceholder{assets/diagrams/image15.png}{15}

Chi spinge in e5 deve esser pronto a rientrare nel Gambetto di Re nel
caso venga giocata 2.e4.

\textbf{2.fxe5 d6 3.exd6 ¥xd6} notare che dopo appena tre mosse il Nero
minaccia già lo scacco matto (4... £h4+ 5.g3 ¥xg3+ 6.hxg3 £xg3 matto).

\textbf{4.¤f3 g5 5.d4 g4}

\emph{Gambetto Jaenisch:} \textbf{1.e4 e5 2.¤f3 ¤c6 3.¥b5 f5}

Questa variante della Partita Spagnola fu ideata dal maestro russo Karl
Jaenisch (1813-1872). Il gambetto sembra giustificato in caso di
accettazione del sacrificio.

\bookdiagramplaceholder{assets/diagrams/image16.png}{16}

\textbf{4.exf5 e4 5.£e2 £e7 6.¥xc6 £xc6 7.¤d4 £e5 8.¤f3 £xf5} il Nero
non esce male dalle schermaglie iniziali però il Bianco può rifiutare
l'offerta con la forte replica 4.¤c3! che porta a situazioni molto
complicate in cui è spesso il Nero a rimetterci le penne.

\emph{Attacco Marshall:} \textbf{1.e4 e5 2.¤f3 ¤c6 3.¥b5 a6 4.¥a4 ¤f6
5.0-0 ¥e7 6.¦e1 b5 7.¥b3 0-0 8.c3 d5}

\bookdiagramplaceholder{assets/diagrams/image17.png}{17}

Il campione americano Frank Marshall tenne per ben nove anni questa idea
``nel cassetto'', in attesa di giocarla in un incontro importante.
L'occasione gliela diede Capablanca a New York nel 1918. Come se nulla
fosse il geniale cubano trovò lì per lì la continuazione ancora oggi
ritenuta migliore e vinse la partita. Il gioco del Nero è stato
successivamente migliorato e oggi questo gambetto è una delle armi più
temibili contro la Partita Spagnola.

\textbf{9.exd5 ¤xd5 10.¤xe5 ¤xe5 11.¦xe5 c6} l'idea originaria di
Marshall fu 11... ¤f6 ma dopo 12.¦e1 ¥d6 13.h3 ¤g4 14.£f3 £h4 15.d4 ¤xf2
16.¦e2! \emph{(16.£xf2? ¥h2+! 17.¢f1 ¥g3. Vale la pena notare la
necessità dello scacco in h2 che obbliga il Re a spostarsi in f1.
All'immediata 16... ¥g3 seguirebbe 17.£xf7+!! ¦xf7 18.¦e8 matto)}
Capablanca respinse l'attacco e vinse la partita.

\textbf{12.d4 ¥d6 13.¦e1 £h4 14.g3 £h3} con buon gioco sulle case
bianche.

\emph{Gambetto Benkö:} \textbf{1.d4 ¤f6 2.c4 c5 3.d5 b5 4.¤xb5 a6ù}

\bookdiagramplaceholder{assets/diagrams/image18.png}{18}

Questa continuazione, nota anche come Gambetto Volga, è del tutto
giocabile ed è l'arma preferita di molti giocatori d'attacco. Il
gambetto ha un sapore del tutto moderno basandosi su un'idea
squisitamente posizionale. Il Nero conquista l'iniziativa sul lato di
Donna facendo pressione sui pedoni a2 e b2 anche grazie al controllo
della diagonale nera a1-h8 e, come se non bastasse, fa pressione sulla
diagonale a6-f1.

\textbf{5.bxa6 ¥xa6 6.¤c3 d6 7.g3 g6 8.¥g2 ¥g7 9.¤f3 0-0 10.0-0 ¤bd7
11.£c2 £a5 12.¦b1 ¦fb8 13.¥d2 ¤b6}

Per finire mettiamo a confronto due modi di giocare lo stesso gambetto.

Le due partite che seguono sono state giocate a distanza di oltre un
secolo, la prima in pieno periodo Romantico, la seconda quando era già
operante la supremazia della scuola sovietica; lo studioso faccia molta
attenzione al diverso spirito che le anima.

Anderssen--Kieseritzky, Londra 1851.

Questa partita, giocata a Londra il 21 giugno 1851, durante una pausa
del primo torneo internazionale della storia degli scacchi, è nota col
nome di ``Immortale''. Essa illustra pe¢fettamente il modo di intendere
gli scacchi all'epoca.

\textbf{1.e4 e5 2.f4 exf4 3.¥c4} in pieno periodo Romantico non si
conoscevano le finezze del gioco posizionale. Ogni mossa o era d'attacco
o era di difesa. In questo ordine di idee si era persa per strada la
concezione di Philidor che i pedoni fossero ``l'anima degli scacchi'';
dal match Bourdonnais-Macdonnell (1834) in poi erano diventati piuttosto
carne da macello. Come è facile immaginare ne scaturivano complicate
baruffe che facevano la gioia degli amanti delle sensazioni forti. Oggi,
si veda la partita successiva, si preferisce la più cauta 3. ¤f3.

\textbf{3... £h4+} all'arrembaggio! Un giocatore moderno opterebbe per
la più semplice (e sana!) 3... ¤f6 oppure per 3... d5.

\textbf{4.¢f1 b5} la mossa preferita da Kieseritsky con la quale si
complica ancora di più il gioco.

\textbf{5.¥xb5 ¤f6 6.¤f3 £h6 7.d3 ch5} minacciando 7... ¤g3+. Si noti in
quanto poco conto siano tenuti i canoni dello sviluppo e si tenda a
creare invece continue minacce con i pochi pezzi già sviluppati.

\textbf{8.ch4 £g5 9.¤f5 c6 10.g4 ¤f6 11.¦g1!} una mossa molto difficile
da trovare. Il Bianco sacrifica l'Alfiere in b5 per guadagnare spazio
sul lato di Re.

\textbf{11... ¤xb5 12.h4 £g6 13.h5 £g5 14.£f3 ¤g8 15.¥xf4 £f6 16.¤c3}
anche un moderno giocatore avrebbe giocato così ma Anderssen la effettua
perché è una mossa d'attacco (minaccia il salto in d5) non per
sviluppare un pezzo.

\textbf{16... ¥c5 17.¤d5 £xb2 18.¥d6!}

\bookdiagramplaceholder{assets/diagrams/image19.png}{19}

Con la precisione di un ragno il Bianco costruisce la sua tela di matto.
Questa mossa implica il sacrificio delle Torri. Adolf Anderssen fu un
maestro insuperato nell'arte della combinazione.

\textbf{18... ¥xg1 19. e5 £xa1+ 20. ¢e2 ¤a6 21. ¤xg7+ ¢d8 22. £f6+!!}
l'o¢gia di sacrifici culmina con quello di Donna.

\textbf{22... ¤xf6 23.¥e7 matto.}

Spassky--Fischer, Mar del Plata 1960.

Quella che segue è una celebre partita, persa dall'allora diciassettenne
Bobby Fischer contro un Boris Spassky quasi all'apice della carriera. Di
scena è ancora il gambetto di Re, giocato però in versione moderna.
All'attacco frontale sul Re qui si lotta per vantaggi di o¢dine
posizionale.

\textbf{1.e4 e5 2.f4 exf4} invece di accettare il pedone il Nero può
impiantare un controgambetto con 2... d5 (Controgambetto Falkbeer)
3.exd5 e4 4.d3 ¤f6 5.¤bd2 exd3 6.¥xd3 ¤xd5. Poiché il secondo tratto del
Bianco non minaccia realmente il \textbf{§}e5 (a 3.fxe5? seguirebbe 3...
£h4+!) il Nero può anche continuare con una semplice mossa di sviluppo:
2... ¥c5 3.¤f3 d6 4.¤c3 ¤f6 5.¥c4 ¤c3 ecc.

\textbf{3.¤f3 g5 4.h4 g4 5.¤e5 ¤f6 6.d4 d6 7.¤d3 ¤xe4 8.¥xf4 ¥g7 9.¤c3?}
le semplificazioni favoriscono il difendente.

\textbf{9... ¤xc3 10.bxc3 c5 11.¥e2 ¤xd4 12.0-0 ¤c6} quella di Fischer
non è né una mossa di difesa né una di attacco ma una semplice mossa di
sviluppo; ai tempi di Anderssen si sarebbe probabilmente preferito una
mossa difensiva quale 12... h5 ma, come fa notare lo stesso Fischer
nelle note alla partita, dopo 12... h5 13.¥g5 f6 14.¥c1 seguita da
15.¤f4 il lato di Re del Nero è tutto buchi. Nel 1858 questa sensibilità
posizionale era di là da venire.

\textbf{13.¥xg4 0-0 14.¥xc8 ¦xc8 15.£g4 f5 16.£g3 £xc3 17.¦ae1 ¢h8
18.¢h1 ¦g8 19.¥xd6 ¥f8 20.¥e5+ ¤xe5 21.£xe5+ ¦g7 22.¦xf5 £xh4+ 23.¢g1
£g4?} occorreva giocare 23... £g3! e dopo 24.£xg3 ¦xg3 il Nero rimane
con un pedone in più.

\textbf{24.¦f2 ¥e7 25.¦e4 £g5 26.£d4 ¦f8?} \emph{``Non vedendo la vera
minaccia del Bianco''} scriverà poi Fischer. \emph{``Io ero preoccupato
per 27.¤e5, senza accorgermi che tale mossa era controbattuta
tranquillamente da 27... ¥c5''.}

\textbf{27.¦e5!} \emph{``Io invece contavo su 27.¤e5? ¦xf2 28.£xf2 ¥c5!
29.£xc5 £xg2 matto.''}

\textbf{27... ¦d8 28.£e4 £h4} \emph{``Sapevo che stavo per perdere un
pezzo, ma ancora non volevo crederci. Dovevo giocare ancora una mossa
per vedere se era proprio vero!''}

\textbf{29.¦f4} il Nero abbandona

\bookdiagramplaceholder{assets/diagrams/image20.png}{20}

L'Alfiere in e7 cade: 29... £h6 30.¦xe7 ¦xe7 \emph{(30... ¦xd3 31.¦f8+
¦g8 32.£e5+ e poi matto)} 31.£xe7 ¦xd3 32.¦f8+ e matto alla successiva.
Notare come il Bianco abbia conquistato per intero il centro.
